\section{Introduction}
\label{sec:intro}

\PARstart{A}{request} such as ``raise the fill volume on Line~1 to 330\,mL'' carries obligations that a
fluent answer does not discharge. The agent must identify the filler and its writable
setpoint rather than a similarly described measurement, interpret the unit, obtain current
evidence, and respect the machine's operating state and process interlocks. These are
supervisory decisions above the controller's deterministic loop; an agent's plausible
response is not evidence that the proposed command is admissible.

Industrial automation has standardised much of the required knowledge over three decades:
\opcua{} information models \cite{iec62541}, the Asset Administration Shell (AAS)
\cite{idta2023aas}, property dictionaries \cite{eclass,iec61360}, ISA-95 hierarchies
\cite{iec62264}, PackML state machines \cite{isa88} and the Module Type Package
\cite{vdi2658}. The semantic web contributes RDF, OWL, SHACL and SPARQL
\cite{w3c2014rdf,w3c2012owl,knublauch2017shacl,w3c2013sparql}, together with QUDT units
\cite{qudt} and SOSA/SSN observations \cite{haller2019ssn}. These technologies were designed
for machine-to-machine interoperability. Whether they supply what a language agent needs at
decision time, and in which combination, has not been analysed systematically.

Existing reviews examine industrial language models, industrial knowledge graphs,
semantic interoperability and individual standards largely independently, and evaluate
technologies against interoperability criteria such as coverage, mapping effort and
conformance. Grounding research, conversely, evaluates retrieval and tool use on
open-domain data that has no units, interlocks or operational consequences
\cite{lewis2020rag,edge2024graphrag,pan2024unifying,patil2024gorilla}. The gap is an
analysis of the field from the perspective of the \emph{engineering evidence} an agent
requires to ground observations and actions.

The central question of this survey is therefore: \emph{what engineering information and
semantic capabilities are required to ground trustworthy industrial AI agents, how well do
existing industrial standards and semantic technologies provide them, and which gaps remain
for deployment?} We contribute:
\begin{enumerate}[label=\textbf{C\arabic*.},leftmargin=2.3em]
  \item A reproducible review protocol with six research questions, five search families,
  explicit eligibility, quality and coding rules, and released screening tooling
  (Section~\ref{sec:method}).
  \item The agent-usability framework (AUF): ten capabilities, each with a definition,
  the agent failure it prevents, a scoring rubric and supporting sources, used as the
  survey's coding instrument (Section~\ref{sec:auf}).
  \item A cross-technology capability analysis of \AufTechnologies{} standards and
  vocabularies organised in a five-layer engineering-semantics landscape
  (Sections~\ref{sec:landscape}--\ref{sec:capabilities}).
  \item A classification of industrial agent architectures by grounding source and
  authority, coded with AUF (Section~\ref{sec:agents}), and a reference synthesis
  derived from it (Section~\ref{sec:architecture}).
  \item A controlled case study showing how capability omissions identified by AUF manifest
  as grounding and validation failures, including a matched \opcua{}/AAS comparison
  (Section~\ref{sec:case}), and a research agenda (Section~\ref{sec:challenges}).
\end{enumerate}
The novelty lies in the analysis perspective, not in any individual technology: we
organise the field around the evidence contract that connects engineering information to
agent failure modes, and map it simultaneously across standards, semantic technologies and
agent architectures.
